\documentclass[12pt,a4paper]{article}

% Packages
\usepackage[utf8]{inputenc}
\usepackage[T1]{fontenc}
\usepackage[a4paper, margin=2.5cm]{geometry}
\usepackage{amsmath}
\usepackage{graphicx}
\usepackage{booktabs}
\usepackage{xcolor}
\usepackage{hyperref}
\usepackage{listings}
\usepackage{caption}
\usepackage{parskip}
\usepackage{enumitem}

% Hyperlink style
\hypersetup{colorlinks=true, linkcolor=blue, urlcolor=blue, citecolor=blue}

% Code listing style
\lstset{
  language=Python,
  basicstyle=\ttfamily\small,
  keywordstyle=\color{blue},
  commentstyle=\color{gray},
  stringstyle=\color{teal},
  breaklines=true,
  frame=single,
  backgroundcolor=\color{gray!10},
  numbers=left,
  numberstyle=\tiny\color{gray},
}

% Title
\title{
  \textbf{Midterm Project} \\
  Real-Time Data Collection and Stream Processing System \\
  \large AITU Classroom Environmental \& Energy Monitoring
}
\author{
  Galymzhan Tamirlan Galymzhanuly \\
  Student ID: 240542 | Group: BDA-2407 \\
  Astana IT University \\
  Instructor: Serek Azamat
}
\date{October 2026}

% ============================================================
\begin{document}
\maketitle
\tableofcontents
\newpage

% ============================================================
\section{Introduction}

For this midterm I decided to build a real-time stream processing system
for monitoring a classroom at AITU. The idea came from Assignment~2, where
we worked with a small batch of classroom sensor data from Room~C1.3.101.
I wanted to take that dataset further and simulate what it would look like
if those sensors were streaming data continuously during a whole lecture.

The goal of the system is to process incoming sensor packets one by one,
compute real-time statistics like rolling averages, and automatically
raise alerts when something unusual happens — bad air quality,
the room getting too hot, or an electrical surge.

I tried to make everything reproducible, so the stream generator uses a fixed
random seed (42) and always produces the same 60-minute sequence.

% ============================================================
\section{Data Scenario and Stream Architecture}

\subsection{The Room and Sensor Setup}

Room~C1.3.101 is a lecture and lab space at AITU with an official capacity
of 40 students. The room has four types of sensors:
\begin{itemize}
  \item \textbf{Temperature sensor} (Temp\_C) — measures indoor air temperature in °C
  \item \textbf{CO\textsubscript{2} sensor} (CO2\_ppm) — measures air quality in parts per million
  \item \textbf{Occupancy camera} (Occupancy) — counts heads using computer vision
  \item \textbf{Smart power meter} (Power\_kW) — measures total electrical load
\end{itemize}

Each sensor sends one reading per minute, timestamped at the minute boundary.
In a real deployment these would go through an MQTT broker or WebSocket
connection. In my simulation I use a Python generator that yields them one at a time.

\subsection{Data Schema}

\begin{table}[h!]
\centering
\caption{Incoming sensor packet schema}
\begin{tabular}{llll}
\toprule
\textbf{Field} & \textbf{Type} & \textbf{Units} & \textbf{Normal Range} \\
\midrule
Time       & string  & HH:MM  & 10:00 – 10:59 \\
Room       & string  & —      & C1.3.101 \\
Temp\_C    & float   & °C     & 20.0 – 24.0 \\
CO2\_ppm   & float   & ppm    & 400 – 1000 \\
Occupancy  & integer & persons & 0 – 40 \\
Power\_kW  & float   & kW     & 1.5 – 7.0 \\
\bottomrule
\end{tabular}
\end{table}

\subsection{Stream Phases}

I structured the 60-minute stream in phases so that several interesting
anomalies appear naturally:

\begin{enumerate}
  \item \textbf{10:00–10:11} — Assignment 2 baseline (12 real records, includes 3 sensor dropouts)
  \item \textbf{10:12–10:25} — Lecture at full capacity, stable readings
  \item \textbf{10:26–10:35} — Ventilation failure, CO\textsubscript{2} spikes past 1000 ppm
  \item \textbf{10:36–10:45} — Thermal accumulation, temperature crosses 24°C
  \item \textbf{10:46–10:50} — Lab practical, electrical surge above 6 kW
  \item \textbf{10:51–10:59} — Lecture ends, students leave, everything drops
\end{enumerate}

% ============================================================
\section{Implementation}

\subsection{Stream Generator (\texttt{stream\_generator.py})}

I wrote a \texttt{generate\_stream()} function that starts with the
Assignment~2 baseline records and then continues with simulated data
for the remaining 48 minutes. The function takes a \texttt{seed} parameter
so the output is always the same.

The three intentional sensor dropouts (missing values) appear at:
\begin{itemize}
  \item 10:02 — Temp\_C missing
  \item 10:04 — CO2\_ppm missing
  \item 10:07 — Occupancy missing
\end{itemize}
Plus three more dropouts I inject in the extended part
(at minutes 18, 33, and 44).

\subsection{Stream Processor (\texttt{stream\_processor.py})}

The \texttt{StreamProcessor} class is the core of the project.
It has a \texttt{process(record)} method that I call for each incoming packet.
Inside, it does the following steps:

\subsubsection{Step 1: Online Imputation}

Missing values are handled before any calculation. I use two strategies:
\begin{itemize}
  \item For continuous sensors (Temp\_C, CO2\_ppm): linear extrapolation
    from the last two buffered values if available, otherwise forward-fill.
  \item For discrete sensors (Occupancy): always forward-fill from the last
    known headcount.
\end{itemize}
Any record with a missing value gets a \texttt{SENSOR\_DROPOUT} flag.

\subsubsection{Step 2: Rolling Statistics}

I keep a sliding window (size 5) for each sensor using Python's
\texttt{deque(maxlen=5)}. For each record I compute:
\begin{align*}
  \bar{T}_t &= \frac{1}{5}\sum_{i=t-4}^{t} T_i
    & \text{(rolling mean temperature)} \\
  \bar{C}_t &= \frac{1}{5}\sum_{i=t-4}^{t} C_i
    & \text{(rolling mean CO\textsubscript{2})} \\
  \bar{P}_t &= \frac{1}{5}\sum_{i=t-4}^{t} P_i
    & \text{(rolling mean power)} \\
  \sigma_T^{(t)} &= \text{std}(T_{t-4} \ldots T_t)
    & \text{(rolling temperature volatility)}
\end{align*}

\subsubsection{Step 3: Derived Indicators}

\begin{itemize}
  \item \textbf{Occupancy density}: $d_t = \text{Occupancy}_t / 40$
  \item \textbf{Energy per occupant}: $e_t = P_t / \max(\text{Occupancy}_t, 1)$
  \item \textbf{Temperature delta}: $\Delta T_t = T_t - T_{t-1}$
  \item \textbf{Comfort flag}: True if $20 \le T_t \le 24$ and $\text{CO}_{2,t} \le 1000$
\end{itemize}

\subsubsection{Step 4: Anomaly Detection}

I check six conditions after each record is processed:

\begin{table}[h!]
\centering
\caption{Anomaly detection rules}
\begin{tabular}{ll}
\toprule
\textbf{Alert} & \textbf{Condition} \\
\midrule
SENSOR\_DROPOUT     & Any sensor field is None (missing) \\
CO2\_HIGH           & CO\textsubscript{2} $> 1000$ ppm \\
CO2\_CRITICAL       & CO\textsubscript{2} $> 1200$ ppm \\
OVERHEATING         & Temp $> 24.0$°C \\
UNDERHEATING        & Temp $< 20.0$°C \\
POWER\_SURGE        & Power $> 6.0$ kW \\
RAPID\_OCC\_CHANGE  & $|\Delta\text{Occupancy}| \ge 6$ persons/min \\
\bottomrule
\end{tabular}
\end{table}

% ============================================================
\section{Results}

After processing all 60 records I got the following statistics:

\begin{itemize}
  \item \textbf{96 anomaly events} detected across 60 records
  \item \textbf{30 OVERHEATING} events (10:36–10:59 mostly)
  \item \textbf{28 POWER\_SURGE} events (10:46–10:59)
  \item \textbf{20 CO2\_CRITICAL} events (10:32–10:45 range)
  \item \textbf{10 CO2\_HIGH} events (10:26–10:31)
  \item \textbf{6 SENSOR\_DROPOUT} events
  \item \textbf{2 RAPID\_OCC\_CHANGE} events (when students leave)
\end{itemize}

The most critical situation was around 10:32–10:35, when the room had
simultaneously high CO\textsubscript{2} (critical level), temperature above 24°C,
and power near the surge threshold.

\subsection{Dashboard Figure}

The figure generated by my notebook (\texttt{stream\_dashboard.png}) shows
three subplots sharing a common time axis:
\begin{enumerate}
  \item Temperature and CO\textsubscript{2} raw values with 5-step rolling means and threshold lines
  \item Power consumption and occupancy over time
  \item Anomaly timeline showing when each alert type fired
\end{enumerate}

% ============================================================
\section{Discussion}

\subsection{What Worked Well}

The online imputation worked better than I expected. Because I use linear
extrapolation from the rolling buffer rather than just forward-filling,
the imputed values follow the trend of the data pretty closely.
For example, when the temperature sensor dropped out at 10:02, the imputed
value (21.8°C) was very close to what would have been recorded.

The anomaly detection is also pretty clean. Because I check multiple conditions
per record and log each one separately, I can see clearly in the timeline
how different alert types overlap — for example, during the ventilation failure
event you can see CO2\_HIGH turning into CO2\_CRITICAL as the concentration builds up.

\subsection{Limitations}

The main limitation is that my stream is simulated. In a real deployment
you'd want to connect to actual sensor hardware. Also, my anomaly thresholds
are static — a smarter system would learn the normal range adaptively for
each room.

Another thing I didn't have time for is sending alerts somewhere useful,
like a notification to the building management system. Right now it all just
gets logged to a DataFrame.

\subsection{What I Would Do Differently}

If I had more time I would add:
\begin{itemize}
  \item Adaptive thresholds using the rolling standard deviation
  \item A proper logging system with timestamps and severity levels
  \item A simple web dashboard using Flask to show the live stream
\end{itemize}

% ============================================================
\section{Conclusion}

I built a working real-time stream processing system that takes classroom
sensor data one record at a time and produces live rolling statistics,
derived comfort indicators, and rule-based alerts.

The key thing I learned is how different stream processing is from batch processing.
In batch mode I can sort, fill, and aggregate the entire dataset at once. In stream
mode I can only use what I've already seen, which forces you to think carefully
about state management and what information to keep in memory.

Using the Assignment~2 dataset as the baseline made it easy to validate the
output — I could compare the first 12 processed records against the values
I computed by hand in Assignment~2 and confirm they matched.

\end{document}
